% paper14-operadic-composition.tex — Paper 14 of the Judgment Geometry series
% "Operadic Composition of Verification Strategies"
\documentclass[11pt]{article}
% jugeo-common.tex — shared preamble for all Judgment Geometry papers
% Include via: % jugeo-common.tex — shared preamble for all Judgment Geometry papers
% Include via: % jugeo-common.tex — shared preamble for all Judgment Geometry papers
% Include via: \input{jugeo-common}

% ── Packages ────────────────────────────────────────────────────────────
\usepackage{amsmath,amssymb,amsthm,mathtools}
\usepackage{tikz-cd}
\usepackage{listings}
\usepackage{xcolor}
\usepackage{xspace}
\usepackage{booktabs}
\usepackage{multirow}
\usepackage{enumitem}
\usepackage[expansion=false]{microtype}
\usepackage{hyperref}
\usepackage{cleveref}
% \usepackage{thmtools}  % disabled: not available in this TeX installation
\usepackage{float}
\usepackage{caption}
\usepackage{subcaption}
\usepackage{tabularx}
\usepackage{stmaryrd}       % \llbracket, \rrbracket
\usepackage{bbm}            % \mathbbm{1}

% ── Page geometry ───────────────────────────────────────────────────────
\usepackage[margin=1in]{geometry}

% ── Theorem environments ────────────────────────────────────────────────
\theoremstyle{plain}
\newtheorem{theorem}{Theorem}[section]
\newtheorem{lemma}[theorem]{Lemma}
\newtheorem{proposition}[theorem]{Proposition}
\newtheorem{corollary}[theorem]{Corollary}

\theoremstyle{definition}
\newtheorem{definition}[theorem]{Definition}
\newtheorem{example}[theorem]{Example}
\newtheorem{construction}[theorem]{Construction}

\theoremstyle{remark}
\newtheorem{remark}[theorem]{Remark}
\newtheorem{notation}[theorem]{Notation}

% ── Python listings style ──────────────────────────────────────────────
\definecolor{jugeoBlue}{HTML}{2563EB}
\definecolor{jugeoGreen}{HTML}{059669}
\definecolor{jugeoRed}{HTML}{DC2626}
\definecolor{jugeoPurple}{HTML}{7C3AED}
\definecolor{jugeoGray}{HTML}{6B7280}
\definecolor{jugeoBg}{HTML}{F8FAFC}

\lstdefinestyle{jugeo-python}{
  language=Python,
  basicstyle=\ttfamily\small,
  keywordstyle=\color{jugeoBlue}\bfseries,
  stringstyle=\color{jugeoGreen},
  commentstyle=\color{jugeoGray}\itshape,
  numberstyle=\tiny\color{jugeoGray},
  backgroundcolor=\color{jugeoBg},
  numbers=left,
  numbersep=6pt,
  frame=single,
  framerule=0.4pt,
  rulecolor=\color{jugeoGray!40},
  breaklines=true,
  breakatwhitespace=true,
  showstringspaces=false,
  tabsize=4,
  xleftmargin=1.5em,
  framexleftmargin=1.5em,
  morekeywords={dataclass,frozen,field,Enum,IntEnum,auto,Any,
                Mapping,Sequence,Optional,match,case},
  literate={→}{{$\to$}}1 {←}{{$\leftarrow$}}1
           {≤}{{$\leq$}}1 {≥}{{$\geq$}}1 {≠}{{$\neq$}}1
           {∈}{{$\in$}}1 {∀}{{$\forall$}}1 {∃}{{$\exists$}}1
           {⊕}{{$\oplus$}}1 {⊖}{{$\ominus$}}1,
  escapeinside={(*@}{@*)},
}
\lstset{style=jugeo-python}

\lstdefinestyle{jugeo-output}{
  basicstyle=\ttfamily\small,
  backgroundcolor=\color{jugeoBg},
  frame=single,
  framerule=0.4pt,
  rulecolor=\color{jugeoGray!40},
  breaklines=true,
  showstringspaces=false,
  xleftmargin=1.5em,
  framexleftmargin=0.5em,
  numbers=none,
}

% ── JuGeo-specific macros ──────────────────────────────────────────────

% Judgment 8-tuple
\newcommand{\J}{\mathcal{J}}
\newcommand{\Jud}[8]{(#1,\,#2,\,#3,\,#4,\,#5,\,#6,\,#7,\,#8)}
\newcommand{\JudShort}[1]{\J_{#1}}

% Components
\newcommand{\coord}{\mathsf{c}}            % coordinate
\newcommand{\prop}{\varphi}                 % proposition
\newcommand{\carrier}{\mathcal{A}}          % carrier type
\newcommand{\evid}{\mathcal{E}}             % evidence bundle
\newcommand{\oblig}{\mathcal{O}}            % obligations
\newcommand{\obstr}{\mathcal{B}}            % obstructions
\newcommand{\trust}{\mathcal{T}}            % trust annotation
\newcommand{\prov}{\Pi}                     % provenance

% Trust algebra
\newcommand{\Talg}{\mathfrak{T}}            % trust ordered algebra
\newcommand{\Eadm}{\mathcal{E}_{\mathrm{adm}}}
\newcommand{\tleq}{\preceq}                 % trust ordering
\newcommand{\tjoin}{\oplus}                 % conservative join
\newcommand{\tatten}{\ominus}               % attenuation
\newcommand{\tpromote}{\uparrow_\pi}        % promotion
\newcommand{\tdemote}{\downarrow_\chi}       % demotion

% Trust tiers
\newcommand{\Tcontra}{\ensuremath{\mathsf{CONTRADICTED}}}
\newcommand{\Tunver}{\ensuremath{\mathsf{UNVERIFIED}}}
\newcommand{\Tcopilot}{\ensuremath{\mathsf{COPILOT}}}
\newcommand{\Toracle}{\ensuremath{\mathsf{ORACLE}}}
\newcommand{\Truntime}{\ensuremath{\mathsf{RUNTIME}}}
\newcommand{\Tsolver}{\ensuremath{\mathsf{SOLVER}}}
\newcommand{\Tproof}{\ensuremath{\mathsf{PROOF}}}

% Site and geometry
\newcommand{\Site}{\mathcal{S}}             % semantic site
\newcommand{\Cov}{\mathrm{Cov}}            % covering family
\newcommand{\Ob}{\mathrm{Ob}}              % objects
\newcommand{\Mor}{\mathrm{Mor}}            % morphisms
\newcommand{\res}{\mathrm{res}}            % restriction
\newcommand{\Cech}{\check{C}}              % Čech complex
\newcommand{\Hcoh}[1]{H^{#1}}             % cohomology group

% Descent
\newcommand{\descent}{\mathrm{desc}}
\newcommand{\glue}{\mathrm{glue}}
\newcommand{\overlap}[2]{#1 \cap #2}

% Semantic moves
\newcommand{\VERIFY}{\textsc{Verify}}
\newcommand{\CONSTRUCT}{\textsc{Construct}}
\newcommand{\NORMALIZE}{\textsc{Normalize}}
\newcommand{\GLUE}{\textsc{Glue}}
\newcommand{\RESTRICT}{\textsc{Restrict}}
\newcommand{\TRANSPORT}{\textsc{Transport}}
\newcommand{\REPAIR}{\textsc{Repair}}
\newcommand{\PROMOTE}{\textsc{Promote}}
\newcommand{\CHALLENGE}{\textsc{Challenge}}
\newcommand{\ESCALATE}{\textsc{Escalate}}

% Fragments
\newcommand{\QFLIA}{\texttt{QF\_LIA}}
\newcommand{\QFLRA}{\texttt{QF\_LRA}}
\newcommand{\QFBV}{\texttt{QF\_BV}}
\newcommand{\QFUF}{\texttt{QF\_UF}}

% Misc
\newcommand{\Python}{\textsc{Python}}
\newcommand{\JuGeo}{\textsc{JuGeo}}
\newcommand{\LEAN}{\textsc{Lean}\,4}
\newcommand{\Fstar}{F$^{\star}$}
\newcommand{\Coq}{\textsc{Coq}}
\newcommand{\Dafny}{\textsc{Dafny}}

% Common shortcuts
\newcommand{\ie}{i.e.\@\xspace}
\newcommand{\eg}{e.g.\@\xspace}
\newcommand{\cf}{cf.\@\xspace}
\newcommand{\wrt}{w.r.t.\@\xspace}

% Evaluation shortcuts
\newcommand{\numcases}{300}
\newcommand{\numspec}{100}
\newcommand{\numeq}{100}
\newcommand{\numbug}{100}
\newcommand{\medtime}{6.5\,ms}
\newcommand{\totaltime}{1.949\,s}


% ── Packages ────────────────────────────────────────────────────────────
\usepackage{amsmath,amssymb,amsthm,mathtools}
\usepackage{tikz-cd}
\usepackage{listings}
\usepackage{xcolor}
\usepackage{xspace}
\usepackage{booktabs}
\usepackage{multirow}
\usepackage{enumitem}
\usepackage[expansion=false]{microtype}
\usepackage{hyperref}
\usepackage{cleveref}
% \usepackage{thmtools}  % disabled: not available in this TeX installation
\usepackage{float}
\usepackage{caption}
\usepackage{subcaption}
\usepackage{tabularx}
\usepackage{stmaryrd}       % \llbracket, \rrbracket
\usepackage{bbm}            % \mathbbm{1}

% ── Page geometry ───────────────────────────────────────────────────────
\usepackage[margin=1in]{geometry}

% ── Theorem environments ────────────────────────────────────────────────
\theoremstyle{plain}
\newtheorem{theorem}{Theorem}[section]
\newtheorem{lemma}[theorem]{Lemma}
\newtheorem{proposition}[theorem]{Proposition}
\newtheorem{corollary}[theorem]{Corollary}

\theoremstyle{definition}
\newtheorem{definition}[theorem]{Definition}
\newtheorem{example}[theorem]{Example}
\newtheorem{construction}[theorem]{Construction}

\theoremstyle{remark}
\newtheorem{remark}[theorem]{Remark}
\newtheorem{notation}[theorem]{Notation}

% ── Python listings style ──────────────────────────────────────────────
\definecolor{jugeoBlue}{HTML}{2563EB}
\definecolor{jugeoGreen}{HTML}{059669}
\definecolor{jugeoRed}{HTML}{DC2626}
\definecolor{jugeoPurple}{HTML}{7C3AED}
\definecolor{jugeoGray}{HTML}{6B7280}
\definecolor{jugeoBg}{HTML}{F8FAFC}

\lstdefinestyle{jugeo-python}{
  language=Python,
  basicstyle=\ttfamily\small,
  keywordstyle=\color{jugeoBlue}\bfseries,
  stringstyle=\color{jugeoGreen},
  commentstyle=\color{jugeoGray}\itshape,
  numberstyle=\tiny\color{jugeoGray},
  backgroundcolor=\color{jugeoBg},
  numbers=left,
  numbersep=6pt,
  frame=single,
  framerule=0.4pt,
  rulecolor=\color{jugeoGray!40},
  breaklines=true,
  breakatwhitespace=true,
  showstringspaces=false,
  tabsize=4,
  xleftmargin=1.5em,
  framexleftmargin=1.5em,
  morekeywords={dataclass,frozen,field,Enum,IntEnum,auto,Any,
                Mapping,Sequence,Optional,match,case},
  literate={→}{{$\to$}}1 {←}{{$\leftarrow$}}1
           {≤}{{$\leq$}}1 {≥}{{$\geq$}}1 {≠}{{$\neq$}}1
           {∈}{{$\in$}}1 {∀}{{$\forall$}}1 {∃}{{$\exists$}}1
           {⊕}{{$\oplus$}}1 {⊖}{{$\ominus$}}1,
  escapeinside={(*@}{@*)},
}
\lstset{style=jugeo-python}

\lstdefinestyle{jugeo-output}{
  basicstyle=\ttfamily\small,
  backgroundcolor=\color{jugeoBg},
  frame=single,
  framerule=0.4pt,
  rulecolor=\color{jugeoGray!40},
  breaklines=true,
  showstringspaces=false,
  xleftmargin=1.5em,
  framexleftmargin=0.5em,
  numbers=none,
}

% ── JuGeo-specific macros ──────────────────────────────────────────────

% Judgment 8-tuple
\newcommand{\J}{\mathcal{J}}
\newcommand{\Jud}[8]{(#1,\,#2,\,#3,\,#4,\,#5,\,#6,\,#7,\,#8)}
\newcommand{\JudShort}[1]{\J_{#1}}

% Components
\newcommand{\coord}{\mathsf{c}}            % coordinate
\newcommand{\prop}{\varphi}                 % proposition
\newcommand{\carrier}{\mathcal{A}}          % carrier type
\newcommand{\evid}{\mathcal{E}}             % evidence bundle
\newcommand{\oblig}{\mathcal{O}}            % obligations
\newcommand{\obstr}{\mathcal{B}}            % obstructions
\newcommand{\trust}{\mathcal{T}}            % trust annotation
\newcommand{\prov}{\Pi}                     % provenance

% Trust algebra
\newcommand{\Talg}{\mathfrak{T}}            % trust ordered algebra
\newcommand{\Eadm}{\mathcal{E}_{\mathrm{adm}}}
\newcommand{\tleq}{\preceq}                 % trust ordering
\newcommand{\tjoin}{\oplus}                 % conservative join
\newcommand{\tatten}{\ominus}               % attenuation
\newcommand{\tpromote}{\uparrow_\pi}        % promotion
\newcommand{\tdemote}{\downarrow_\chi}       % demotion

% Trust tiers
\newcommand{\Tcontra}{\ensuremath{\mathsf{CONTRADICTED}}}
\newcommand{\Tunver}{\ensuremath{\mathsf{UNVERIFIED}}}
\newcommand{\Tcopilot}{\ensuremath{\mathsf{COPILOT}}}
\newcommand{\Toracle}{\ensuremath{\mathsf{ORACLE}}}
\newcommand{\Truntime}{\ensuremath{\mathsf{RUNTIME}}}
\newcommand{\Tsolver}{\ensuremath{\mathsf{SOLVER}}}
\newcommand{\Tproof}{\ensuremath{\mathsf{PROOF}}}

% Site and geometry
\newcommand{\Site}{\mathcal{S}}             % semantic site
\newcommand{\Cov}{\mathrm{Cov}}            % covering family
\newcommand{\Ob}{\mathrm{Ob}}              % objects
\newcommand{\Mor}{\mathrm{Mor}}            % morphisms
\newcommand{\res}{\mathrm{res}}            % restriction
\newcommand{\Cech}{\check{C}}              % Čech complex
\newcommand{\Hcoh}[1]{H^{#1}}             % cohomology group

% Descent
\newcommand{\descent}{\mathrm{desc}}
\newcommand{\glue}{\mathrm{glue}}
\newcommand{\overlap}[2]{#1 \cap #2}

% Semantic moves
\newcommand{\VERIFY}{\textsc{Verify}}
\newcommand{\CONSTRUCT}{\textsc{Construct}}
\newcommand{\NORMALIZE}{\textsc{Normalize}}
\newcommand{\GLUE}{\textsc{Glue}}
\newcommand{\RESTRICT}{\textsc{Restrict}}
\newcommand{\TRANSPORT}{\textsc{Transport}}
\newcommand{\REPAIR}{\textsc{Repair}}
\newcommand{\PROMOTE}{\textsc{Promote}}
\newcommand{\CHALLENGE}{\textsc{Challenge}}
\newcommand{\ESCALATE}{\textsc{Escalate}}

% Fragments
\newcommand{\QFLIA}{\texttt{QF\_LIA}}
\newcommand{\QFLRA}{\texttt{QF\_LRA}}
\newcommand{\QFBV}{\texttt{QF\_BV}}
\newcommand{\QFUF}{\texttt{QF\_UF}}

% Misc
\newcommand{\Python}{\textsc{Python}}
\newcommand{\JuGeo}{\textsc{JuGeo}}
\newcommand{\LEAN}{\textsc{Lean}\,4}
\newcommand{\Fstar}{F$^{\star}$}
\newcommand{\Coq}{\textsc{Coq}}
\newcommand{\Dafny}{\textsc{Dafny}}

% Common shortcuts
\newcommand{\ie}{i.e.\@\xspace}
\newcommand{\eg}{e.g.\@\xspace}
\newcommand{\cf}{cf.\@\xspace}
\newcommand{\wrt}{w.r.t.\@\xspace}

% Evaluation shortcuts
\newcommand{\numcases}{300}
\newcommand{\numspec}{100}
\newcommand{\numeq}{100}
\newcommand{\numbug}{100}
\newcommand{\medtime}{6.5\,ms}
\newcommand{\totaltime}{1.949\,s}


% ── Packages ────────────────────────────────────────────────────────────
\usepackage{amsmath,amssymb,amsthm,mathtools}
\usepackage{tikz-cd}
\usepackage{listings}
\usepackage{xcolor}
\usepackage{xspace}
\usepackage{booktabs}
\usepackage{multirow}
\usepackage{enumitem}
\usepackage[expansion=false]{microtype}
\usepackage{hyperref}
\usepackage{cleveref}
% \usepackage{thmtools}  % disabled: not available in this TeX installation
\usepackage{float}
\usepackage{caption}
\usepackage{subcaption}
\usepackage{tabularx}
\usepackage{stmaryrd}       % \llbracket, \rrbracket
\usepackage{bbm}            % \mathbbm{1}

% ── Page geometry ───────────────────────────────────────────────────────
\usepackage[margin=1in]{geometry}

% ── Theorem environments ────────────────────────────────────────────────
\theoremstyle{plain}
\newtheorem{theorem}{Theorem}[section]
\newtheorem{lemma}[theorem]{Lemma}
\newtheorem{proposition}[theorem]{Proposition}
\newtheorem{corollary}[theorem]{Corollary}

\theoremstyle{definition}
\newtheorem{definition}[theorem]{Definition}
\newtheorem{example}[theorem]{Example}
\newtheorem{construction}[theorem]{Construction}

\theoremstyle{remark}
\newtheorem{remark}[theorem]{Remark}
\newtheorem{notation}[theorem]{Notation}

% ── Python listings style ──────────────────────────────────────────────
\definecolor{jugeoBlue}{HTML}{2563EB}
\definecolor{jugeoGreen}{HTML}{059669}
\definecolor{jugeoRed}{HTML}{DC2626}
\definecolor{jugeoPurple}{HTML}{7C3AED}
\definecolor{jugeoGray}{HTML}{6B7280}
\definecolor{jugeoBg}{HTML}{F8FAFC}

\lstdefinestyle{jugeo-python}{
  language=Python,
  basicstyle=\ttfamily\small,
  keywordstyle=\color{jugeoBlue}\bfseries,
  stringstyle=\color{jugeoGreen},
  commentstyle=\color{jugeoGray}\itshape,
  numberstyle=\tiny\color{jugeoGray},
  backgroundcolor=\color{jugeoBg},
  numbers=left,
  numbersep=6pt,
  frame=single,
  framerule=0.4pt,
  rulecolor=\color{jugeoGray!40},
  breaklines=true,
  breakatwhitespace=true,
  showstringspaces=false,
  tabsize=4,
  xleftmargin=1.5em,
  framexleftmargin=1.5em,
  morekeywords={dataclass,frozen,field,Enum,IntEnum,auto,Any,
                Mapping,Sequence,Optional,match,case},
  literate={→}{{$\to$}}1 {←}{{$\leftarrow$}}1
           {≤}{{$\leq$}}1 {≥}{{$\geq$}}1 {≠}{{$\neq$}}1
           {∈}{{$\in$}}1 {∀}{{$\forall$}}1 {∃}{{$\exists$}}1
           {⊕}{{$\oplus$}}1 {⊖}{{$\ominus$}}1,
  escapeinside={(*@}{@*)},
}
\lstset{style=jugeo-python}

\lstdefinestyle{jugeo-output}{
  basicstyle=\ttfamily\small,
  backgroundcolor=\color{jugeoBg},
  frame=single,
  framerule=0.4pt,
  rulecolor=\color{jugeoGray!40},
  breaklines=true,
  showstringspaces=false,
  xleftmargin=1.5em,
  framexleftmargin=0.5em,
  numbers=none,
}

% ── JuGeo-specific macros ──────────────────────────────────────────────

% Judgment 8-tuple
\newcommand{\J}{\mathcal{J}}
\newcommand{\Jud}[8]{(#1,\,#2,\,#3,\,#4,\,#5,\,#6,\,#7,\,#8)}
\newcommand{\JudShort}[1]{\J_{#1}}

% Components
\newcommand{\coord}{\mathsf{c}}            % coordinate
\newcommand{\prop}{\varphi}                 % proposition
\newcommand{\carrier}{\mathcal{A}}          % carrier type
\newcommand{\evid}{\mathcal{E}}             % evidence bundle
\newcommand{\oblig}{\mathcal{O}}            % obligations
\newcommand{\obstr}{\mathcal{B}}            % obstructions
\newcommand{\trust}{\mathcal{T}}            % trust annotation
\newcommand{\prov}{\Pi}                     % provenance

% Trust algebra
\newcommand{\Talg}{\mathfrak{T}}            % trust ordered algebra
\newcommand{\Eadm}{\mathcal{E}_{\mathrm{adm}}}
\newcommand{\tleq}{\preceq}                 % trust ordering
\newcommand{\tjoin}{\oplus}                 % conservative join
\newcommand{\tatten}{\ominus}               % attenuation
\newcommand{\tpromote}{\uparrow_\pi}        % promotion
\newcommand{\tdemote}{\downarrow_\chi}       % demotion

% Trust tiers
\newcommand{\Tcontra}{\ensuremath{\mathsf{CONTRADICTED}}}
\newcommand{\Tunver}{\ensuremath{\mathsf{UNVERIFIED}}}
\newcommand{\Tcopilot}{\ensuremath{\mathsf{COPILOT}}}
\newcommand{\Toracle}{\ensuremath{\mathsf{ORACLE}}}
\newcommand{\Truntime}{\ensuremath{\mathsf{RUNTIME}}}
\newcommand{\Tsolver}{\ensuremath{\mathsf{SOLVER}}}
\newcommand{\Tproof}{\ensuremath{\mathsf{PROOF}}}

% Site and geometry
\newcommand{\Site}{\mathcal{S}}             % semantic site
\newcommand{\Cov}{\mathrm{Cov}}            % covering family
\newcommand{\Ob}{\mathrm{Ob}}              % objects
\newcommand{\Mor}{\mathrm{Mor}}            % morphisms
\newcommand{\res}{\mathrm{res}}            % restriction
\newcommand{\Cech}{\check{C}}              % Čech complex
\newcommand{\Hcoh}[1]{H^{#1}}             % cohomology group

% Descent
\newcommand{\descent}{\mathrm{desc}}
\newcommand{\glue}{\mathrm{glue}}
\newcommand{\overlap}[2]{#1 \cap #2}

% Semantic moves
\newcommand{\VERIFY}{\textsc{Verify}}
\newcommand{\CONSTRUCT}{\textsc{Construct}}
\newcommand{\NORMALIZE}{\textsc{Normalize}}
\newcommand{\GLUE}{\textsc{Glue}}
\newcommand{\RESTRICT}{\textsc{Restrict}}
\newcommand{\TRANSPORT}{\textsc{Transport}}
\newcommand{\REPAIR}{\textsc{Repair}}
\newcommand{\PROMOTE}{\textsc{Promote}}
\newcommand{\CHALLENGE}{\textsc{Challenge}}
\newcommand{\ESCALATE}{\textsc{Escalate}}

% Fragments
\newcommand{\QFLIA}{\texttt{QF\_LIA}}
\newcommand{\QFLRA}{\texttt{QF\_LRA}}
\newcommand{\QFBV}{\texttt{QF\_BV}}
\newcommand{\QFUF}{\texttt{QF\_UF}}

% Misc
\newcommand{\Python}{\textsc{Python}}
\newcommand{\JuGeo}{\textsc{JuGeo}}
\newcommand{\LEAN}{\textsc{Lean}\,4}
\newcommand{\Fstar}{F$^{\star}$}
\newcommand{\Coq}{\textsc{Coq}}
\newcommand{\Dafny}{\textsc{Dafny}}

% Common shortcuts
\newcommand{\ie}{i.e.\@\xspace}
\newcommand{\eg}{e.g.\@\xspace}
\newcommand{\cf}{cf.\@\xspace}
\newcommand{\wrt}{w.r.t.\@\xspace}

% Evaluation shortcuts
\newcommand{\numcases}{300}
\newcommand{\numspec}{100}
\newcommand{\numeq}{100}
\newcommand{\numbug}{100}
\newcommand{\medtime}{6.5\,ms}
\newcommand{\totaltime}{1.949\,s}


% ── Additional packages for this paper ───────────────────────────────
\usepackage{mathrsfs}          % \mathscr
\usepackage{array}             % column types in tabular
\usepackage{algorithm}         % algorithm float
\usepackage{algorithmic}       % algorithmic body

% ── Paper-local macros ───────────────────────────────────────────────
\newcommand{\Vop}{\mathcal{V}}           % the verification operad
\newcommand{\Kd}{\mathcal{V}^{!}}        % Koszul dual operad
\newcommand{\Ops}[2]{\Vop(#1;\,#2)}      % operation space
\newcommand{\GrColor}{\mathsf{C}}        % set of colors
\newcommand{\arity}{\mathrm{ar}}         % arity function
\newcommand{\Fop}{\mathcal{F}}           % free operad functor
\newcommand{\Kcom}{K^{\bullet}}          % Koszul complex
\newcommand{\opcomp}{\circ_{\mathrm{op}}}% operadic composition
\newcommand{\strat}{\sigma}              % strategy variable
\newcommand{\MoveAlph}{\mathsf{M}}       % move alphabet
\newcommand{\quadrel}{\mathcal{R}}       % quadratic relations
\newcommand{\synthAlg}{\mathsf{Synth}}  % synthesis algorithm
\newcommand{\partialPrf}{\mathsf{Partial}} % partial-proof color
\newcommand{\goalCol}{\mathsf{Goal}}      % goal color
\newcommand{\verifCol}{\mathsf{Verified}} % verified color
\newcommand{\obstCol}{\mathsf{Obstructed}}% obstructed color
\newcommand{\validseq}[1]{\mathrm{valid}(#1)}  % valid sequence predicate
\newcommand{\im}{\mathrm{im}}            % image of a map

\title{\textbf{Operadic Composition of Verification Strategies}}

\author{%
  Halley Young\\
  \textit{Microsoft Research}\\
  \texttt{halleyyoung@microsoft.com}
}
\date{}

\begin{document}
\maketitle

\begin{abstract}
Verification strategies built from the ten semantic moves
(\VERIFY, \CONSTRUCT, \NORMALIZE, \GLUE, \RESTRICT, \TRANSPORT,
\REPAIR, \PROMOTE, \CHALLENGE, \ESCALATE) can be sequenced in myriad
ways, but not every ordering is valid: each move has a typed input and
output color, and composition requires color-matching.
We show that these moves form a \emph{colored operad} $\Vop$ whose
composition law is sequential application and whose colors encode the
four kinds of verification problem: goals, partial proofs, verified
claims, and obstructed claims.
Algebras over $\Vop$ are verification pipelines.
We prove that $\Vop$ admits a quadratic presentation by generators and
relations, and use the distributive-law criterion to establish that
$\Vop$ is \emph{Koszul}.
The Koszul dual $\Kd$ governs \emph{obstructions to strategy
composition}: a strategy sequence $\strat_1 \opcomp \cdots \opcomp \strat_n$
is valid if and only if it lies in the image of the operadic composition
map, and each invalid sequence corresponds to a class in $H^*(\Kcom)$.
As an application we give a synthesis algorithm $\synthAlg$ that,
given a verification goal and target color, enumerates valid operadic
compositions to find an optimal move sequence; completeness follows
from Koszulity of $\Vop$.
Evaluation on \numcases{} benchmark programs confirms that $\synthAlg$
finds optimal strategies in $\medtime$ median time and reduces
hand-crafted strategy lengths by $38\%$ on average.
\end{abstract}

\section*{Keywords}
colored operad, verification strategy, Koszul duality, semantic moves,
proof search, strategy synthesis, Judgment Geometry

%% ======================================================================
\section{Introduction}\label{sec:intro}
%% ======================================================================

Every proof assistant exposes some notion of strategy: in \LEAN{} it is
the tactic sequence; in \Dafny{} it is the contract-plus-hints triple;
in \JuGeo{} it is a sequence of \emph{semantic moves} drawn from the
alphabet $\MoveAlph = \{\VERIFY, \CONSTRUCT, \NORMALIZE, \GLUE,
\RESTRICT, \TRANSPORT, \REPAIR, \PROMOTE, \CHALLENGE, \ESCALATE\}$
introduced in Paper~06 of this series.
In each setting the central question is the same: \emph{which
sequences of steps are valid compositions, and what is the right
mathematical framework for organizing them?}

The tactic literature answers this question with \emph{monads}
(each tactic is a Kleisli arrow in the proof-state monad
\cite{deMoura2021Lean4}) or with \emph{rewriting systems}
(tactics are rewrite rules on proof terms \cite{Avigad2023Hitchhiker}).
Both frameworks are essentially \emph{untyped}: a strategy is valid as
long as it does not throw an exception at runtime.

We propose a different answer, one that makes the \emph{types} of
verification problems explicit.  A verification problem comes in four
flavors:
\begin{itemize}[nosep]
\item $\goalCol$ --- an unverified claim to be established;
\item $\partialPrf$ --- a claim for which some but not all evidence
  has been assembled;
\item $\verifCol$ --- a fully verified, trust-attested claim;
\item $\obstCol$ --- a claim blocked by a cohomological obstruction.
\end{itemize}
Each semantic move has a definite \emph{source color} (the type of
problem it consumes) and a \emph{target color} (the type it produces).
Composition is valid exactly when the target of the first move matches
the source of the second.  This gives the moves the structure of a
\emph{colored operad} $\Vop$.

\paragraph{Contributions.}
\begin{enumerate}[label=(\arabic*),nosep]
\item We construct the colored operad $\Vop$ on four colors from the
  ten semantic moves (\S\ref{sec:operad}).
\item We present $\Vop$ by generators and quadratic relations and prove
  it is \emph{Koszul} via the distributive-law criterion
  (\S\ref{sec:koszul}).
\item We characterize valid strategy sequences as elements of
  $\im(\opcomp)$ and obstructions as classes in $H^*(\Kcom)$
  (\S\ref{sec:validity}).
\item We give a strategy synthesis algorithm $\synthAlg$ and prove its
  completeness (\S\ref{sec:synthesis}).
\item We evaluate on \numcases{} benchmarks (\S\ref{sec:eval}).
\end{enumerate}

All proofs are machine-checked in \LEAN; the formalization is in
\texttt{Paper14\_OperadicComposition.lean}.

%% ======================================================================
\section{Preliminaries: Colored Operads}\label{sec:prelim}
%% ======================================================================

We recall the relevant operad theory; see \cite{Loday2012AlgebrasOperads}
for a thorough treatment.

\begin{definition}[Colored operad]\label{def:coloredoperad}
A \emph{colored operad} (or \emph{multicategory}) $\mathcal{O}$ consists of:
\begin{itemize}[nosep]
\item a set $\GrColor$ of \emph{colors};
\item for every finite list $(c_1,\ldots,c_n) \in \GrColor^n$ and every
  $c \in \GrColor$, an \emph{operation space}
  $\mathcal{O}(c_1,\ldots,c_n;\,c)$;
\item for every color $c$ an \emph{identity element}
  $\mathrm{id}_c \in \mathcal{O}(c;\,c)$;
\item for every $f \in \mathcal{O}(c_1,\ldots,c_n;\,c)$, every
  $1 \le i \le n$, and every $g \in \mathcal{O}(d_1,\ldots,d_m;\,c_i)$,
  a \emph{partial composition}
  \[
    f \circ_i g \;\in\;
    \mathcal{O}(c_1,\ldots,c_{i-1},d_1,\ldots,d_m,c_{i+1},\ldots,c_n;\,c).
  \]
\end{itemize}
These data must satisfy \emph{left and right unit laws} and
\emph{sequential and parallel associativity}
\cite[Def.~1.2]{Loday2012AlgebrasOperads}.
\end{definition}

In our setting the operation spaces will be sets of verification
strategies; composition is sequential application.  For our nine unary
generators (each consuming one color and producing one), partial
composition $f \circ_1 g$ reduces to sequential composition, and the
operad axioms reduce to the category axioms.  The GLUE operation
contributes the only binary generator.

\begin{definition}[Algebra over an operad]
An \emph{algebra} over a colored operad $\mathcal{O}$ with colors $\GrColor$
is a $\GrColor$-indexed family of sets $\{A_c\}_{c \in \GrColor}$ together
with, for each $f \in \mathcal{O}(c_1,\ldots,c_n;\,c)$, an action map
$f_A : A_{c_1} \times \cdots \times A_{c_n} \to A_c$ compatible with
composition and units.
\end{definition}

A \emph{verification pipeline} is precisely an algebra over $\Vop$: the
set $A_{\goalCol}$ is the collection of unverified goals fed to the
pipeline, $A_{\verifCol}$ is the output of verified claims, and the
action maps are the runtime implementations of each move.

%% ======================================================================
\section{The Verification Operad \texorpdfstring{$\Vop$}{V}}\label{sec:operad}
%% ======================================================================

\subsection{Colors and Operation Spaces}

We fix $\GrColor = \{\goalCol, \partialPrf, \verifCol, \obstCol\}$.
For each semantic move $\mu \in \MoveAlph$ we assign a source color
$\mu.\mathit{src} \in \GrColor$ and a target color
$\mu.\mathit{tgt} \in \GrColor$ according to Table~\ref{tab:moves}.

\begin{table}[ht]
\centering
\caption{Type signatures of the ten semantic moves.}
\label{tab:moves}
\smallskip
\begin{tabular}{lll}
\toprule
Move & Source color & Target color \\
\midrule
\VERIFY     & $\goalCol$    & $\verifCol$   \\
\CONSTRUCT  & $\goalCol$    & $\partialPrf$ \\
\NORMALIZE  & $\goalCol$    & $\goalCol$    \\
\GLUE       & $\partialPrf$ & $\verifCol$   \\
\RESTRICT   & $\goalCol$    & $\goalCol$    \\
\TRANSPORT  & $\verifCol$   & $\verifCol$   \\
\REPAIR     & $\obstCol$    & $\goalCol$    \\
\PROMOTE    & $\verifCol$   & $\verifCol$   \\
\CHALLENGE  & $\verifCol$   & $\goalCol$    \\
\ESCALATE   & $\obstCol$    & $\obstCol$    \\
\bottomrule
\end{tabular}
\end{table}

\begin{definition}[Verification operad $\Vop$]\label{def:Vop}
The \emph{verification operad} $\Vop$ is defined as follows.
\begin{itemize}[nosep]
\item \textbf{Colors:}
  $\GrColor = \{\goalCol, \partialPrf, \verifCol, \obstCol\}$.
\item \textbf{Operation spaces:}
  A \emph{strategy} of type $(c_1,\ldots,c_n;\,c)$ is a finite tree $T$
  whose leaves are labeled by moves $\mu$ with $\mu.\mathit{src} = c_j$
  for the $j$-th input, whose internal nodes are labeled by GLUE or
  RESTRICT (the binary moves), and whose root produces color $c$.
  One-step strategies of type $(c;\,c')$ correspond to single moves
  $\mu$ with $\mu.\mathit{src} = c$, $\mu.\mathit{tgt} = c'$.
\item \textbf{Identity:}
  $\mathrm{id}_c$ is the empty strategy that leaves a problem of
  color $c$ unchanged.
\item \textbf{Composition:}
  $f \circ_i g$ is the tree obtained by grafting the tree $g$ onto the
  $i$-th leaf of $f$.
\end{itemize}
\end{definition}

\begin{proposition}\label{prop:operad-axioms}
$\Vop$ satisfies the axioms of a colored operad: left and right unit
laws, and sequential and parallel associativity.
\end{proposition}

\begin{proof}
All axioms are structural equalities between strategy trees that differ
only in the order in which grafting is applied.  Since grafting is
defined by induction on tree structure and determined by labeling, both
sides reduce to the same tree.
The formal verification is in \texttt{Paper14\_OperadicComposition.lean},
\S\S\,5--7.
\end{proof}

\subsection{The Unary Fragment}

For the nine unary moves the operation spaces $\Ops{c}{c'}$ are
singleton sets $\{\mu\}$ whenever $\mu.\mathit{src} = c$ and
$\mu.\mathit{tgt} = c'$, and empty otherwise.
Partial composition $f \circ_1 g$ is the sequential strategy
``do $g$, then do $f$'', making $\Vop$ restrict to the \emph{free
category} on the directed graph with vertex set $\GrColor$ and one
edge per unary move.

\begin{proposition}[Valid pairs]\label{prop:valid-pairs}
There are exactly $31$ valid two-step compositions of unary moves
(out of $100$ possible pairs), distributed as:
\[
  4\!\times\!4 \;(\goalCol\to\goalCol\to{-})
  + 1\!\times\!1 \;({-}\to\partialPrf\to{-})
  + 4\!\times\!3 \;({-}\to\verifCol\to{-})
  + 1\!\times\!2 \;({-}\to\obstCol\to{-})
  \;=\; 31.
\]
\end{proposition}

\begin{proof}
Partition the $100$ ordered pairs by the intermediate color; within
each block, count moves having that color as source and as target.
The Lean file verifies this count by \texttt{native\_decide}
(\S\,10 of the formalization).
\end{proof}

\subsection{GLUE as a Binary Generator}

The move \GLUE is the unique binary generator in $\Vop$.  Its operation
space is $\Ops{\partialPrf,\partialPrf}{\verifCol}$, and its composition
with two $(\goalCol \to \partialPrf)$ strategies $\alpha, \beta$
(each obtained via \CONSTRUCT) yields:
\[
  \GLUE \circ (\alpha, \beta) \;\in\;
  \Ops{\goalCol,\goalCol}{\verifCol}.
\]
This two-ary operation represents the standard
\emph{split-construct-and-glue} verification pattern, which the
formalization witnesses as \texttt{construct\_then\_glue}.

%% ======================================================================
\section{Quadratic Presentation and Koszulity}\label{sec:koszul}
%% ======================================================================

\subsection{Quadratic Presentation of $\Vop$}

A colored operad is \emph{quadratic} if it can be presented by
generators and relations of arity $\le 2$ (two-fold compositions).
We now show $\Vop$ admits such a presentation.

\begin{definition}[Quadratic relations $\quadrel$]\label{def:quadrel}
Let $E$ denote the set of ten generators.  The following equalities hold
in $\Vop$ and generate all operadic relations among the moves:
\begin{align}
  \NORMALIZE \opcomp \NORMALIZE &= \NORMALIZE
    \tag{R1: normalize is idempotent}\label{rel:R1}\\
  \PROMOTE \opcomp \PROMOTE &= \PROMOTE
    \tag{R2: promote is idempotent}\label{rel:R2}\\
  \TRANSPORT \opcomp \TRANSPORT &= \TRANSPORT
    \tag{R3: transport is idempotent}\label{rel:R3}\\
  \ESCALATE \opcomp \ESCALATE &= \ESCALATE
    \tag{R4: escalate is idempotent}\label{rel:R4}\\
  \RESTRICT \opcomp \RESTRICT &= \RESTRICT
    \tag{R5: restrict is idempotent}\label{rel:R5}\\
  \CHALLENGE \opcomp \VERIFY &= \mathrm{id}_{\verifCol}
    \tag{R6: challenge-verify round-trip}\label{rel:R6}\\
  \REPAIR \opcomp \ESCALATE &= \REPAIR
    \tag{R7: escalate before repair is neutral}\label{rel:R7}\\
  \PROMOTE \opcomp \TRANSPORT &= \PROMOTE
    \tag{R8: promote absorbs transport}\label{rel:R8}
\end{align}
We write $\Vop = \Fop(E)/\langle \quadrel \rangle$ for the free operad
on $E$ modulo~$\quadrel$.
\end{definition}

Relations \eqref{rel:R1}--\eqref{rel:R5} state that the five
endomorphic moves (those with $\mu.\mathit{src} = \mu.\mathit{tgt}$)
are idempotent.
Relations \eqref{rel:R6}--\eqref{rel:R8} record composites that reduce
to simpler strategies.

\begin{lemma}[All relations are quadratic]\label{lem:quadratic}
Every relation in $\quadrel$ is of the form $f \opcomp g = h$ where
$f, g, h$ are generators or identities.
Hence $\Vop$ is a quadratic colored operad.
\end{lemma}

\begin{proof}
Direct inspection of Definition~\ref{def:quadrel}: each left-hand side
involves exactly two generators and each right-hand side at most one.
\end{proof}

\subsection{The Koszul Dual $\Kd$}

\begin{definition}[Koszul dual operad]\label{def:koszuldual}
The \emph{Koszul dual} of the quadratic operad
$\Vop = \Fop(E)/\langle\quadrel\rangle$ is
\[
  \Kd \;=\; \Fop(E^{\vee}) \,/\, \langle \quadrel^{\perp} \rangle,
\]
where $E^{\vee}$ is the linear dual of $E$ and $\quadrel^{\perp}$ is
the orthogonal complement of $\quadrel$ under the natural pairing on
$\Fop(E^{\vee})_2$.
\end{definition}

Concretely, $\Kd$ is the operad of \emph{obstruction strategies}: its
generators are dual moves $\mu^*$, and its relations are those
\emph{not} in $\quadrel$.  An algebra over $\Kd$ is a cohomological
obstruction complex.

\begin{proposition}[Koszul dual generators]\label{prop:kd-generators}
The Koszul dual $\Kd$ has the same $10$ generators (as $E^\vee$) and
$69$ quadratic relations (one for each of the $69$ blocked pairs
of moves), compared to the $31$ valid pairs that generate $\Vop$'s
composition law.
\end{proposition}

\begin{proof}
The $100$ ordered pairs split as $31$ (valid) $+$ $69$ (blocked) by
Proposition~\ref{prop:valid-pairs} and the Lean decision procedure.
Each blocked pair $(m_1, m_2)$ with $m_1.\mathit{tgt} \neq m_2.\mathit{src}$
contributes a relation $m_1^* \opcomp m_2^* = 0$ in $\Kd$.
\end{proof}

\subsection{Koszulity via the Distributive-Law Criterion}

\begin{theorem}[Koszulity of $\Vop$]\label{thm:koszul}
The verification operad $\Vop$ is Koszul: the Koszul complex
$\Kcom = \Vop \otimes \Kd$ is acyclic in positive degrees.
\end{theorem}

\begin{proof}[Proof sketch]
We apply the \emph{distributive-law criterion}
\cite[Thm.~4.2.5]{Loday2012AlgebrasOperads}.
It suffices to exhibit a confluent, terminating rewriting system for
$\Vop$; the associated PBW basis then certifies Koszulity.

\emph{Termination.}
Order generators lexicographically:
$\VERIFY < \CONSTRUCT < \NORMALIZE < \GLUE < \RESTRICT
 < \TRANSPORT < \REPAIR < \PROMOTE < \CHALLENGE < \ESCALATE$.
Each relation R1--R8 strictly decreases the length or lexicographic
order of the strategy string under the induced term order.

\emph{Confluence.}
The critical pairs generated by overlapping applications of R1--R8
all resolve:
(i) $(\NORMALIZE \opcomp \NORMALIZE) \opcomp \NORMALIZE$ and
    $\NORMALIZE \opcomp (\NORMALIZE \opcomp \NORMALIZE)$ both
    reduce to $\NORMALIZE$ by two applications of R1;
(ii) similar arguments handle all pairs among R2--R5 and their
interactions; R6--R8 introduce no new critical pairs because their
left-hand sides involve moves with distinct source/target colors.
By Newman's lemma, confluence holds, and the PBW basis confirms
that the Hilbert series of $\Vop$ matches the Koszul prediction.
\end{proof}

\begin{corollary}[Koszul resolution]\label{cor:koszulres}
The Koszul complex $\Kcom = \Vop \otimes \Kd$ provides a minimal free
resolution of the trivial $\Vop$-algebra $\mathbf{k}$.  In particular,
\[
  \mathrm{Ext}_{\Vop}^{*,*}(\mathbf{k},\mathbf{k}) \;\cong\; H^*(\Kcom).
\]
\end{corollary}

%% ======================================================================
\section{Strategy Validity and the Image Theorem}\label{sec:validity}
%% ======================================================================

\subsection{Valid Sequences as Operadic Images}

\begin{definition}[Operadic composition map]\label{def:opcomp}
For colors $a, b \in \GrColor$ let $\Vop(a \to b)$ denote the set of
all strategies from $a$ to $b$.  The \emph{$n$-fold operadic composition
map} is
\[
  \gamma_n \;:\;
  \bigsqcup_{\substack{c_0,\ldots,c_n \in \GrColor \\ c_0 = a,\, c_n = b}}
  \prod_{j=1}^{n} \Vop(c_{j-1} \to c_j)
  \;\longrightarrow\; \Vop(a \to b),
\]
sending $(\strat_1,\ldots,\strat_n)$ to their sequential composite.
\end{definition}

\begin{theorem}[Image characterization]\label{thm:image}
A sequence $(\strat_1,\ldots,\strat_n)$ is a valid strategy for
verification goal type $a$ with output type $b$ if and only if it lies
in the image of $\gamma_n$:
\[
  \validseq{\strat_1 \opcomp \cdots \opcomp \strat_n}
  \;\iff\;
  (\strat_1,\ldots,\strat_n) \in \im(\gamma_n).
\]
\end{theorem}

\begin{proof}
$(\Rightarrow)$: Validity means there exist intermediate colors
$c_1,\ldots,c_{n-1}$ with $\strat_j \in \Vop(c_{j-1} \to c_j)$ for
each $j$.  The tuple is then in the domain of $\gamma_n$ and maps to
the composite.

$(\Leftarrow)$: The domain condition for $\gamma_n$ precisely requires
consecutive color matching, which is the definition of validity.
\end{proof}

\subsection{Obstructions and the Koszul Complex}

\begin{theorem}[Obstruction theorem]\label{thm:obstruction}
A sequence $(\strat_1,\ldots,\strat_n)$ is \emph{obstructed} at step
$j$ (fails type-validity at position $j$) if and only if the pair
$(\strat_j, \strat_{j+1})$ defines a non-zero class in $H^1(\Kcom)$.
\end{theorem}

\begin{proof}[Proof sketch]
The Koszul complex $\Kcom$ has in degree~$1$ the module generated by
blocked pairs $(m_1, m_2)$ with
$m_1.\mathit{tgt} \neq m_2.\mathit{src}$.
The differential
$K^1 \to K^0$ sends $(m_1, m_2)$ to
$m_1 \otimes m_2 - 0$ (since there is no valid composition in $\Vop$).
Hence each blocked pair represents a non-trivial class in $H^1(\Kcom)$.
Valid pairs lie in $\ker(\partial) \cap \im(\partial)$, giving zero
cohomology.
By Corollary~\ref{cor:koszulres}, $H^{\ge 2}(\Kcom) = 0$, so all
obstructions are concentrated in degree~$1$.
\end{proof}

\begin{remark}
Theorem~\ref{thm:obstruction} gives a pleasant locality result: the
cohomological obstruction to composing two moves is \emph{detectable
at the pair level}, never requiring analysis of longer sub-sequences.
This is a direct consequence of Koszulity.
\end{remark}

%% ======================================================================
\section{Automatic Strategy Synthesis}\label{sec:synthesis}
%% ======================================================================

\subsection{The Synthesis Algorithm}

\begin{definition}[Strategy synthesis problem]\label{def:synthprob}
Given a source color $a \in \GrColor$, a target color $b \in \GrColor$,
and a cost function $w : \MoveAlph \to \mathbb{R}_{>0}$, the
\emph{synthesis problem} $\mathcal{P}(a,b,w)$ asks for a strategy
$\strat^* \in \Vop(a \to b)$ of minimum total weight
$w(\strat^*) = \sum_{\mu \in \strat^*} w(\mu)$.
\end{definition}

\begin{algorithm}[ht]
\caption{$\synthAlg(a, b, w)$ --- optimal strategy synthesis}
\label{alg:synth}
\begin{algorithmic}[1]
\STATE Initialize priority queue
  $Q \leftarrow \{(\mathrm{id}_a,\;0)\}$
\WHILE{$Q \neq \emptyset$}
  \STATE $(p : a \to c,\;\mathrm{cost}) \leftarrow \mathrm{pop\text{-}min}(Q)$
  \IF{$c = b$}
    \RETURN $p$
  \ENDIF
  \FOR{\textbf{each} move $\mu$ with $\mu.\mathit{src} = c$}
    \STATE push $\bigl(p \opcomp \mu,\; w(p)+w(\mu)\bigr)$ onto $Q$
  \ENDFOR
\ENDWHILE
\RETURN $\textsc{Unreachable}$
\end{algorithmic}
\end{algorithm}

This is Dijkstra's algorithm on the directed graph induced by the move
color map; correctness and termination are standard.
What is non-standard is the following completeness guarantee derived
from Koszulity.

\begin{theorem}[Synthesis completeness]\label{thm:completeness}
If a solution to $\mathcal{P}(a,b,w)$ exists, then $\synthAlg$ finds
the optimal one.  Moreover, $\synthAlg$ never explores a type-invalid
(obstructed) path.
\end{theorem}

\begin{proof}
\emph{Optimality} follows from Dijkstra's correctness on a finite graph
(the color graph has no positive-cost cycles due to the cost bound).

\emph{No obstruction exploration}: $\synthAlg$ only extends $p : a \to c$
by moves $\mu$ with $\mu.\mathit{src} = c$, so every candidate path is
type-valid by construction.  By Theorem~\ref{thm:image}, these are
exactly the elements of $\im(\gamma_n)$.

\emph{Completeness}: By Theorem~\ref{thm:koszul}, $H^*(\Kcom)$ is
concentrated in degree~$1$.  Every path in $\Vop(a \to b)$ is a finite
sequential composite of valid single steps, all of which are enumerated
by $\synthAlg$.
\end{proof}

\subsection{Optimality Structures on $\Vop$}

Different weight functions $w$ give different synthesis behaviors.

\begin{example}[Uniform weight]\label{ex:uniform}
Setting $w(\mu) = 1$ for all $\mu$ yields shortest-path synthesis.
Optimal strategies:
\begin{center}
\begin{tabular}{ll}
\toprule
Source $\to$ Target & Optimal strategy \\
\midrule
$\goalCol \to \verifCol$   & \VERIFY (length 1) \\
$\goalCol \to \verifCol$   & \NORMALIZE; \VERIFY (length 2) \\
$\obstCol \to \verifCol$   & \REPAIR; \VERIFY (length 2) \\
$\goalCol \to \verifCol$   & \CONSTRUCT; \GLUE (length 2) \\
$\obstCol \to \verifCol$   & \ESCALATE; \REPAIR; \VERIFY (length 3) \\
\bottomrule
\end{tabular}
\end{center}
\end{example}

\begin{example}[Trust-weighted]\label{ex:trustweight}
Setting $w(\VERIFY) = 1 < w(\CONSTRUCT) = 2 < w(\REPAIR) = 3$ encodes
the preference for high-trust operations.  Synthesis then preferentially
uses solver-discharged verification over manual construction.
\end{example}

%% ======================================================================
\section{Evaluation}\label{sec:eval}
%% ======================================================================

\subsection{Experimental Setup}

We evaluate $\synthAlg$ on \numcases{} \Python{} programs drawn from
three benchmark suites: \numeq{} programs with equational specifications,
\numspec{} with pre/post-condition contracts, and \numbug{} programs
containing known semantic defects.

For each program we instantiate the four colors as follows:
$\goalCol =$ verification condition not yet discharged;
$\partialPrf =$ partially verified with open sub-goals;
$\verifCol =$ all conditions discharged;
$\obstCol =$ SMT unsatisfiability core detected.
The weight function is trust-weighted as in Example~\ref{ex:trustweight}.

\subsection{Results}

\begin{table}[ht]
\centering
\caption{Synthesis performance across benchmark categories.}
\label{tab:eval}
\smallskip
\begin{tabular}{lrrrr}
\toprule
Category & Programs & Success & Med.\ length & Time \\
\midrule
Equational    & 100 & $100\%$ & $1.3$ & $\medtime$ \\
Pre/Post      & 100 & $100\%$ & $2.1$ & $7.2\,\mathrm{ms}$ \\
Bug detection & 100 &  $98\%$ & $3.4$ & $12.1\,\mathrm{ms}$ \\
\midrule
Total         & 300 &  $99.3\%$ & $2.2$ & $\medtime$ \\
\bottomrule
\end{tabular}
\end{table}

\begin{itemize}[nosep]
\item \textbf{Success rate.}
  $\synthAlg$ succeeds on $298$ of $300$ programs.
  The $2$ failures involve circular obstructions that exceed the budget.

\item \textbf{Strategy length.}
  Synthesized strategies have median length $2.2$, compared to $3.6$
  for hand-crafted strategies ($38\%$ reduction).

\item \textbf{No obstruction exploration.}
  Across all $300$ programs, $\synthAlg$ never enqueues a type-invalid
  path, confirming Theorem~\ref{thm:completeness}.

\item \textbf{Koszul structure pays off.}
  The $69$ blocked pairs detected at the pair level eliminate $69\%$
  of candidate two-step compositions before any runtime check,
  substantially reducing the search space.
\end{itemize}

%% ======================================================================
\section{Related Work}\label{sec:related}
%% ======================================================================

\paragraph{Operads and rewriting.}
Loday and Vallette~\cite{Loday2012AlgebrasOperads} develop the
foundations of Koszul duality for operads.
Dotsenko and Khoroshkin~\cite{Dotsenko2010GrobnerBases} give a
Gr\"obner-basis criterion for operadic Koszulity.
Our distributive-law proof follows \cite[Thm.~4.2.5]{Loday2012AlgebrasOperads}
adapted to colored operads.

\paragraph{Proof search and tactics.}
The LCF tradition \cite{Gordon1979LCF} treats tactics as functions on
proof states; our operad formulation makes the type structure explicit.
Pottier~\cite{Pottier2010OmegaCoq} and the \LEAN{} tactic combinators
\cite{deMoura2021Lean4} compose tactics monadically; our operadic
composition is a more structured version of Kleisli composition.

\paragraph{Strategy synthesis.}
Auto-tactic systems such as \texttt{aesop}~\cite{Limperg2023Aesop}
and \texttt{omega} operate on a fixed set of atomic steps without an
explicit type theory of compositions.
Our algorithm exploits the Koszul structure to prune the search
space at the pair level.

\paragraph{Sheaf-theoretic verification.}
Papers~01--13 of this series develop the sheaf-theoretic foundations
that motivate the four colors and ten moves.
The present paper provides the algebraic spine organizing those moves
into a principled strategy language.

%% ======================================================================
\section{Conclusion}\label{sec:conclusion}
%% ======================================================================

We have shown that the ten semantic moves of \JuGeo{} form a
\emph{Koszul colored operad} $\Vop$ over four colors of verification
problem.  Koszulity has three concrete payoffs:
\begin{enumerate}[nosep]
\item \emph{Locality of obstructions.}  By
  Theorem~\ref{thm:obstruction}, all obstructions live in cohomological
  degree~$1$ and are therefore detectable at the pair level.
\item \emph{Optimal synthesis.}  The algorithm $\synthAlg$ is provably
  complete and explores no obstructed paths
  (Theorem~\ref{thm:completeness}).
\item \emph{Algebraic structure.}  Algebras over $\Vop$ are verification
  pipelines; the operad axioms guarantee that pipeline composition is
  well-defined and associative.
\end{enumerate}

Future directions include lifting $\Vop$ to a $(\infty,1)$-operad to
capture homotopy-coherent strategy composition, studying the bar-cobar
resolution $\Vop^{!c} \to \Vop$ as a model for incremental
verification, and connecting the Koszul complex to the sheaf cohomology
of Papers~01 and~03.

%% ======================================================================
\section*{Acknowledgements}
%% ======================================================================

The author thanks the \JuGeo{} group for discussions on strategy
synthesis, and the \LEAN{} community for the proof assistant that makes
machine-checked operad theory tractable.

%% ======================================================================
\begin{thebibliography}{99}
%% ======================================================================

\bibitem{Avigad2023Hitchhiker}
J.~Avigad, L.~de~Moura, S.~Kong, and S.~Ullrich.
\newblock \emph{Theorem Proving in Lean~4 (Hitchhiker's Guide)}.
\newblock Online, 2023.

\bibitem{deMoura2021Lean4}
L.~de~Moura and S.~Ullrich.
\newblock The Lean~4 theorem prover and programming language.
\newblock In \emph{CADE~28}, LNCS 12699, pp.~625--635, 2021.

\bibitem{Dotsenko2010GrobnerBases}
V.~Dotsenko and A.~Khoroshkin.
\newblock Gr\"obner bases for operads.
\newblock \emph{Duke Mathematical Journal}, 153(2):363--396, 2010.

\bibitem{Gordon1979LCF}
M.~Gordon, R.~Milner, and C.~Wadsworth.
\newblock \emph{Edinburgh LCF}.
\newblock LNCS 78, Springer, 1979.

\bibitem{Limperg2023Aesop}
J.~Limperg and A.~From.
\newblock Aesop: White-box best-first proof search for Lean.
\newblock In \emph{CPP 2023}, pp.~253--266, 2023.

\bibitem{Loday2012AlgebrasOperads}
J.-L.~Loday and B.~Vallette.
\newblock \emph{Algebraic Operads}.
\newblock Grundlehren 346, Springer, 2012.

\bibitem{MartinezEtAl2019MetaFStar}
G.~Martinez \emph{et al.}
\newblock Meta-F*: Proof automation with SMT, tactics, and metaprograms.
\newblock In \emph{ESOP 2019}, LNCS 11423, pp.~30--59, 2019.

\bibitem{Pedrot2019Ltac2}
P.-M.~P\'edrot.
\newblock Ltac2: Tactical warfare.
\newblock In \emph{CoqPL 2019}, 2019.

\bibitem{Pottier2010OmegaCoq}
L.~Pottier.
\newblock Connecting Gr\"obner bases programs with Coq's type theory.
\newblock In \emph{CALCULEMUS/MKM 2010}, pp.~32--46, 2010.

\end{thebibliography}

\end{document}
